La separación de fuentes musicales o sonoras es, a día de hoy, un campo abierto en el ámbito de la investigación del procesado de señal. Se puede considerar como fuente sonora toda aquella que genere sonido susceptible de ser captado por un sensor de audio, es decir,
un micrófono. Por tanto, son fuentes sonoras los instrumentos musicales, la voz, o cualquier fuente de ruido natural o artificial.


El problema en sí mismo se atribuye al ámbito del procesamiento digital de señales y en él se plantea el siguiente reto: desarrollar métodos para que, partiendo de una pieza musical, se obtengan los elementos sonoros que la componen, como por ejemplo obtener a partir de una pieza en formato de archivo digital
del género de cantautor donde confluyen una vocal masculina y un instrumento, por ejemplo, una guitarra acústica. El reto consiste en  obtener un método que tras su aplicación nos proporcione por un lado una pista con la información sonora de la vocal, y por otro lado otra con la de la guitarra acústica.  

Los métodos que se desarrollan para dar solución a este problema tienen aplicaciones en distintos ámbitos relacionados con el procesamiento de audio. Entre ellas se encuentran la creación de versiones instrumentales o de karaoke, la remezcla de obras musicales, la restauración de grabaciones, la identificación y extracción de voces, el análisis automático de música y el apoyo a tareas de transcripción o clasificación sonora. También constituye un problema relevante para el estudio de técnicas de aprendizaje profundo aplicadas a señales temporales.

El trabajo propuesto presenta una dificultad considerable debido a que las fuentes pueden compartir regiones temporales y frecuenciales. Por ejemplo, una voz y un instrumento pueden producir energía en frecuencias similares o sonar simultáneamente y como consecuencia no existe una división directa de la pista que permita recuperar cada fuente sin introducir errores, interferencias o ruido de otros elementos no deseados. Por otro lado, también 
existen diferentes técnicas de mezcla y masterización que pueden alterar la señal original en función de cómo interactúan las pistas que componen la canción entre ellas, efectos de reverberación, compresión, ecualización, etc. Estos factores modifican la forma de onda de las diferentes fuentes cuya suma compone la mezcla final y dificultan la generalización, que suele ser el objetivo de los métodos de separación de fuentes.

En el proyecto se ha abordado el problema de la separación de fuentes musicales mediante las herramientas que nos proporcionan los avances tecnológicos de la última década, aplicando inteligencia artificial y aprendizaje automático. La separación de fuentes aún siendo un ámbito muy explorado en el pasado, sigue suponiendo un reto en nuestro tiempo, las pistas pueden compartir regiones temporales y/o frecuenciales, una voz y un instrumento pueden producir energía en frecuencias cercanas del espectro frecuencial y como consecuencia no existir una división directa de ambos que nos permita obtener dos elementos separados. Además, existen factores como algunas técnicas concretas de mezcla o los efectos aplicados de reverberación, compresión, ecualización, etcétera, que hacen que obtener
una generalización  mediante aprendizaje automático se haga una tarea compleja.

El resto de la memoria se aborda de la siguiente manera: en el capítulo 2 se comenta la planificación seguida a lo largo del proyecto y en el capítulo 3 las herramientas que se han utilizado para su desarrollo. Posteriormente, en el capítulo 4, se hace una introducción al marco teórico del procesamiento de señales y del aprendizaje automático, explicando conceptos necesarios para entender el proyecto. En el capítulo 5 se hace un recorrido a lo largo del estado del arte del ámbito. En el capítulo 6 se explican los dos entregables principales que se proporcionan en el trabajo y finalmente en los capítulos siete y ocho, se dan explicaciones técnicas acerca de la
integración del sistema y de su evaluación respectivamente.

Al final del documento en los capítulos nueve y diez se encuentran las conclusiones finales y el trabajo futuro, junto con la sección de anexos. A continuación de esta última sección se encuentra la bibliografía.